\subsection{Compact localization at the base mixed exponents}
\label{subsec:compact-base-cases}
\begin{corollary}[The base quintic and septic arithmetic sources]
\label{cor:compact-base-source-five-seven}
For $a\in\{5,7\}$ and $S_a=\{3,a\}$, the map
\[
 H^1(G_{\Q,S_a},T_3J_a)\longrightarrow H^1(\Q_3,T_3J_a)
\]
is injective integrally. In particular the compact localization
$\Sel(T_3J_a)\to H^1(\Q_3,T_3J_a)$ is injective.
Together with the pure-power and mixed-family source
arguments and Theorem~\ref{thm:eleven-compact-selmer}, this proves the abelian localization assertions of
Theorem~\ref{thm:compact-localization-sections-intro}.
\end{corollary}
\begin{proof}
The quintic action and local-source theorem were proved first on
$U_5$ itself, before passing through the three-power covers. The
same is true at seven, where Proposition~\ref{prop:lm-whole-septic-action}
explicitly includes $r=0$. The verified class-group certificate
supplies its arithmetic-source surjectivity. Thus the cocycle-graph
argument applies to the compact coefficient modules $T_3J_5$ and
$T_3J_7$ directly. More explicitly, a cocycle on $G_{\Q,S_a}$
and its coefficient action define a pro-three graph over the
residual field. Both factor through the established arithmetic
source. Its image from $G_{\Q_3}$ is the entire source; after
subtracting the same integral local coboundary globally, the
cocycle is zero. This proves the full cohomology injection,
including at seven where the residual-field degree is divisible
by three and cannot be used for prime-to-three averaging.

A compact Selmer class is unramified outside $\{3,5\}$ or
$\{3,7\}$, respectively: at any other prime the abelian variety
has good reduction and the geometric Kummer class is unramified.
The integral cohomology localization injections from the source
theorems therefore kill any such class that is zero at three.
This argument uses neither the cubic local detector nor separation
of cusp differences by their three-primary order. Those additional
steps are used only in the affine classification, whose mixed
families still require $r\ge1$.
\end{proof}
